\section{Implementation and Methodology}\label{sec:method}

Every result in this report is produced by running code. This section describes the implementation, how
we checked it, and how we measured. The repository's README lists the command behind each result.

\subsection{Implementation}
Both schemes, all six algorithms each, are implemented in Python~3.12 over native group libraries called
through \code{ctypes}. The reference implementation of the repaired scheme (\code{scheme.py}) is a direct
transliteration of the paper's Section~V-A, with the paper's variable names, so it can be read side by
side with the paper. Optimised variants live in separate modules, so the reference keeps measuring the
algorithm the paper describes. Two printed typos in the paper are corrected in the code: \textsf{KeyGen} prints $pk_i=s_i$ for
$pk_i=s_iP$, and the repaired certificate check uses Verma's $R_{1i}$ where $R_i$ is meant.

\paragraph{Two independent group backends.} \prov{Implementation}
The paper's algebra divides by hash values and its proof relies on a group of prime order. Curve25519
(Ed25519)~\cite{bernstein2012} has cofactor 8, so we use Ristretto255~\cite{rfc9496}, a prime-order group
built on it, through libsodium~1.0.18. A second backend implements NIST P-256~\cite{fips186} through
OpenSSL~3.0.13. The two libraries share no code and represent points differently, so every test is also a
differential test: two independent implementations must agree on every signature, verification result
and operation count. P-256 is also the curve most often deployed in industrial stacks.

\paragraph{Hashing and encoding.} \prov{Implementation}
The paper applies $H_1$ and $H_2$ to identical inputs. Instantiating both with the same hash function
would make $u_i=v_i$ and silently remove the repair's third change. Each oracle therefore gets its own
domain-separation tag, and a test asserts $u_i\neq v_i$ on identical input. Hash inputs are encoded with
length prefixes, because plain concatenation of variable-length fields is ambiguous
(\code{"ab"}$\concat$\code{"c"} and \code{"a"}$\concat$\code{"bc"} collide). Each hash is SHA-512 over the
tagged encoding, reduced modulo the group order from 64 bytes, which makes the bias negligible.

\paragraph{Operation counting.}
Every scalar multiplication, point addition and hash passes through one counting wrapper, so the
$\Te$/$\Ta$/$\Th$ counts of \cref{sec:cost} are measured, not derived by hand. The same instrument
reproduces Verma's published cost rows exactly, which validates the method.

\subsection{Testing}
The test suite has 424 tests: 199 that run once per backend (398) and 26 that need no group, such as the
curve-fitting and aggregator-model tests. A standalone self-test checks 62 algebraic properties of the two
backends before anything else runs. The tests cover round-trip correctness for both schemes up to
$n=17$, rejection of every kind of tampering, the forgeries and their failures, both ablations, the
operation counts, the optimised verifiers against the reference, the aggregator policies, nonce-reuse
recovery and the hedged-nonce fix.

\subsection{Measurement methodology}
Timings are taken on an AMD Ryzen~5 7520U laptop, with a second machine (an Intel Xeon at 2.2~GHz on
Kaggle) used to confirm the timing results. A first sweep on the laptop produced a spurious result,
$n=100$ appearing cheaper than $n=50$, because the CPU clock roughly doubled partway through the sweep.
Two measures address this. A stabilisation step runs real scalar multiplications until the per-operation
cost has stayed constant for several consecutive chunks and at least four seconds have passed. Every
measurement then interleaves all workloads, across both schemes and all values of $n$, takes the median
of several rounds, and calibrates repetitions so that each sample lasts at least 20~ms.

We report \emph{ratios} of per-signer cost, obtained as the slope of a linear fit of aggregate
verification time against $n$ ($n=50$ to $500$), rather than absolute times. Absolute times include Python
and \code{ctypes} overhead and depend on the curve, which is not the paper's: the paper's unit costs come
from a pairing library and it does not name its curve. A ratio between two schemes on the same machine,
curve and library is the quantity the paper's tables make a claim about.

\subsection{Evidence discipline}
Every number in this report and in the accompanying slides is recomputed by one audit script, which
reruns the attacks, counts and tests and reads the recorded benchmark files, asserts each claim, and
stops if any fails. Deterministic results are reported as ``$k/8$'': 4 target messages on each of the 2
backends. These trials confirm that the code behaves as the algebra predicts on both libraries; they are
not statistical samples.
